% 2015 年陕西卷（文史类）第 7 题：程序框图（依固定扫描页重绘，只画拓扑与标号，不按比例）
\begin{tikzpicture}[>=stealth,scale=0.95]
  % 开始
  \draw[rounded corners=6pt] (-1.0,0.4) rectangle (1.0,-0.4);
  \node at (0,0) {开始};
  \draw[->] (0,-0.4)--(0,-1.0);
  % 输入 x
  \draw (-1.0,-1.0)--(1.2,-1.0)--(1.0,-1.8)--(-1.2,-1.8)--cycle;
  \node at (0,-1.4) {输入 $x$};
  \draw[->] (0,-1.8)--(0,-2.3);
  % 赋值 x = x - 3
  \draw (-1.3,-2.3) rectangle (1.3,-3.1);
  \node at (0,-2.7) {$x=x-3$};
  \draw[->] (0,-3.1)--(0,-3.7);
  % 判断 x >= 0
  \draw (-1.9,-4.4)--(0,-3.7)--(1.9,-4.4)--(0,-5.1)--cycle;
  \node at (0,-4.4) {$x\ge 0$};
  % 是：回到赋值框
  \draw[->] (1.9,-4.4)--(3.2,-4.4)--(3.2,-2.7)--(1.3,-2.7);
  \node at (2.35,-4.05) {是};
  % 否：向下
  \draw[->] (0,-5.1)--(0,-5.7);
  \node at (0.32,-5.4) {否};
  % y = x^2 + 1
  \draw (-1.4,-5.7) rectangle (1.4,-6.5);
  \node at (0,-6.1) {$y=x^{2}+1$};
  \draw[->] (0,-6.5)--(0,-7.1);
  % 输出 y
  \draw (-1.0,-7.1)--(1.2,-7.1)--(1.0,-7.9)--(-1.2,-7.9)--cycle;
  \node at (0,-7.5) {输出 $y$};
  \draw[->] (0,-7.9)--(0,-8.5);
  % 结束
  \draw[rounded corners=6pt] (-1.0,-8.5) rectangle (1.0,-9.3);
  \node at (0,-8.9) {结束};
\end{tikzpicture}
